%% SCRAP: architecture/03-architecture/physics-engine/steady-state-machine
%% SOURCE: docs/working/architecture/03-architecture/physics-engine/steady-state-machine.md
%% STATUS: CURRENT
%% FITS: dev-guide/ch-physics
%% EDITORIAL: lifted — prose rewritten to press voice

\section{The Steady-State Machine}

A Steady-State Machine (SSM) is a computational model whose operational state is
defined not by discrete symbolic transitions but by convergence to an optimal
runtime equilibrium. Where a Finite State Machine, pushdown automaton, or Turing
Machine makes ``state'' explicit and symbolic, an SSM derives its behavior from
continuous runtime variables. Using execution statistics as a proxy for
thermodynamic quantities, the SSM tracks execution heat, entropy (workload
diversity), pipeline readiness, heat decay, the steady-state slope, localized
temperature gradients, and adaptive strategy selection. The machine observes
itself, adapts, and settles into stable performance basins called steady states.
In an SSM the attractor defines the machine, not the symbol.

\subsection{Core Definition}

An SSM is a computational system defined by a set of internal metrics, a set of
adaptive policies, an attractor-based state space, and a governing dynamic. The
metrics are execution heat $H$, entropy and diversity measures $E$, pipeline
activation thresholds $P$, a heat-decay profile $\Theta$, and the sliding
execution window $W$. The adaptive policies are lookup strategies, caching
regimes, bucket reorganizations, pipeline decisions, and inference-based mode
switches. The state space is composed of attractor basins --- steady,
quasi-steady, transient, and chaotic --- rather than symbolic states. The
governing dynamic is

\begin{equation}
M_{t+1} = f\bigl(M_t,\; \Delta H,\; \Delta W,\; \Delta E\bigr),
\end{equation}

where $f$ drives the machine toward a stable basin. The machine is defined by its
convergence behavior, not by its instruction sequence.

\subsection{The Fundamental Insight}

In an SSM, performance \emph{is} the state. The system continuously measures how
hot words are, how often patterns recur, how pipeline-able a sequence is, how
diverse the window is, and how much the dictionary should reorganize. As these
quantities stabilize, the machine settles into a self-maintaining optimal regime,
preserved until the environment changes --- through new workloads or a diversity
spike --- at which point the machine passes through a transient or chaotic regime
before finding a new equilibrium.

\subsection{Phases}

An SSM moves through four phases:

\begin{itemize}
  \item \textbf{Transient.} The system warms up; execution patterns are
        sporadic, heat gradients noisy, and the pipeline inactive or unstable.
  \item \textbf{Quasi-steady.} Patterns emerge, lookups begin to bias, hot words
        form local attractors, and dictionary reorders stabilize.
  \item \textbf{True steady state.} The heat surface is smooth, the pipeline
        active, the diversity window predictable, and lookups converge to
        optimal; runtime exceeds the baseline VM and the system remains stable
        unless perturbed.
  \item \textbf{Chaotic/disruption.} Triggered by a major workload shift,
        dictionary mutation, an extreme entropy spike, or a cold-cache shock, the
        SSM falls out of equilibrium temporarily.
\end{itemize}

\subsection{Axioms}

\begin{itemize}
  \item \textbf{Axiom 1 --- Convergence.} Every sustained workload induces the
        SSM to converge toward a stable attractor unless external entropy forces
        divergence.
  \item \textbf{Axiom 2 --- Locality of Heat.} Execution heat reflects both
        locality and temporal relevance; locality produces stability.
  \item \textbf{Axiom 3 --- Adaptive Reordering.} Reordering is not symbolic
        mutation but thermodynamic self-optimization.
  \item \textbf{Axiom 4 --- Stability Maximizes Throughput.} The steady state is
        always faster than the cold state and usually faster than the naive
        baseline VM.
  \item \textbf{Axiom 5 --- Perturbation Response.} When disrupted, the SSM seeks
        a new steady state; it does not thrash indefinitely.
\end{itemize}

\subsection{Relationship to Other Computational Models}

The SSM is distinguished from established models by its emergent, non-symbolic
state. A Finite State Machine occupies discrete symbolic states, whereas the SSM
occupies emergent attractors in a metric space. A Turing Machine defines behavior
through tape symbols and transitions; the SSM defines it through adaptation toward
equilibrium. A Markov chain moves probabilistically between explicit states; the
SSM drifts deterministically toward stable basins under observation. A neural
network learns weights; the SSM reorganizes runtime structures rather than
parameters.

\subsection{Novelty}

The SSM is presented as the first runtime model whose state is emergent rather
than explicit: a runtime self-optimization model, an adaptive VM with convergence
properties, a load-dependent reordering machine, a self-optimizing dictionary
machine, and a phase-shifting execution system. It defines a new category of
adaptive computation.

\subsection{Canonical Example: StarForth}

StarForth is the reference implementation of the SSM. It realizes execution heat,
rolling diversity windows, phase-based lookup strategies, adaptive dictionary
ordering, pipeline activation, a DoE-calibrated inference loop, and a background
physics monitor.

\subsection{Formal Summary}

A Steady-State Machine is a computational system in which the operational state is
defined by attractor convergence in a runtime metric space rather than by
discrete symbolic transitions. An SSM continuously measures workload
characteristics, adjusts its internal structures, and stabilizes into
performance-optimized steady states; when perturbed, it transitions through
transient or chaotic phases before reaching a new equilibrium. This model defines
a class of adaptive computational systems with properties distinct from finite
state machines, pushdown automata, and Turing Machines.
